% chapters/09/section-05.tex
\section{Program Mesin Keadaan Lengkap}
\label{sec:program-fsm}

Bagian ini menyatukan seluruh unsur menjadi program NPC bermesin keadaan yang
dapat dimainkan dan dapat diuji.

\subsection{Kode Lengkap}

\kodepython{python/ch09_fsm.py}{NPC dengan mesin keadaan empat keadaan}{lst:ch09-fsm}

\subsection{Membedah Susunan Mesin Keadaannya}

Seluruh perilaku NPC terletak pada kelas \code{Npc}, dan tersusun sebagai
berikut.

\begin{description}
  \item[\code{decide}] Membaca keadaan sekarang, jarak terhadap pemain, dan
        ketiga pewaktu, lalu mengembalikan nama keadaan berikutnya. Fungsi ini
        tidak mengubah keadaan apa pun.
  \item[\code{enter_state}] Menetapkan keadaan baru, mengembalikan penghitung
        keadaan itu ke nol, dan mencatat perpindahannya.
  \item[\code{act}] Menjalankan perilaku keadaan yang sedang aktif, yaitu
        mengubah kecepatan NPC.
  \item[\code{update}] Menyatukan ketiganya: menghitung jarak, meminta
        keputusan, menyiapkan keadaan baru apabila berpindah, lalu bertindak.
\end{description}

Fungsi \code{update} hanya berisi empat langkah dan tidak memuat satu
pun syarat perpindahan. Seluruh syarat itu berada di dalam \code{decide}.
Susunan itulah yang membuat keputusan dapat diperiksa terpisah dari
perilakunya.

\subsection{Keluaran Pengujian}

Ketika program dijalankan dengan skenario pengujian, terminal menampilkan
catatan perpindahan beserta hasil pemeriksaannya.

\begin{kodekeluaran}
Bingkai dijalankan : 900
Serangan dilakukan : 5
Catatan perpindahan keadaan (bingkai, keadaan):
        0  berdiam
       91  berpatroli
      151  mengejar
      217  menyerang
      421  mengejar
      482  berdiam
      571  berpatroli
Hasil pemeriksaan  : seluruh perpindahan sah
Hasil skenario     : SESUAI dengan urutan yang diharapkan
\end{kodekeluaran}

Ketujuh perpindahan pada Tabel~\ref{tab:tujuh-perpindahan} terjadi seluruhnya,
dalam urutan yang dapat diperkirakan, dan seluruhnya sah menurut tabel
perpindahan. Empat keadaan yang diperlihatkan pada
Gambar~\ref{fig:diagram-keadaan} telah diuji satu per satu.

\subsection{Percobaan yang Disarankan}

\begin{table}[htbp]
\centering
\caption{Percobaan pada program mesin keadaan}
\label{tab:percobaan-ch09}
{\setlength{\tabcolsep}{4pt}\footnotesize%
\begin{tabularx}{\textwidth}{@{}L{5.6cm} Y@{}}
\toprule
\headrow \textbf{Yang diubah dan caranya} & \textbf{Yang sebaiknya Anda amati} \\
\midrule
Pewaktu jeda serangan \code{ATTACK_COOLDOWN}, misalnya menjadi 0,2 detik. &
Jumlah serangan pada catatan bertambah, sedangkan urutan keadaannya tidak
berubah. Hitunglah jumlah yang seharusnya dan bandingkan dengan catatan. \\
\addlinespace
Pewaktu tenggang \code{GIVE_UP_SECONDS}, misalnya menjadi 3 detik. &
NPC mengejar lebih lama sebelum menyerah. Pemain harus berdiri
lebih jauh agar NPC benar-benar kehilangan jejak. \\
\addlinespace
Pewaktu berdiam \code{IDLE_SECONDS}, misalnya menjadi nol. &
Keadaan berdiam tidak lagi berlangsung lama; NPC langsung berpatroli sesudah
menyerah. \\
\addlinespace
Penghalang jangkauan serang, dengan menghapus perkalian 1,2. &
NPC berkedip dan tersendat pada tepi jangkauan serang, sebab kedua syarat
perpindahan benar bersamaan. \\
\addlinespace
Sebuah syarat perpindahan, dengan menambahkan perpindahan yang tidak ada pada
diagram, misalnya dari menyerang langsung ke berpatroli. &
Pemeriksaan keabsahan pada terminal melaporkan perpindahan yang tidak sah. \\
\bottomrule
\end{tabularx}}
\end{table}

\begin{peringatanbox}
Percobaan terakhir adalah yang paling penting untuk dicoba. Menambahkan satu
perpindahan yang tidak ada pada diagram akan langsung ditangkap oleh
pemeriksaan keabsahan. Itulah manfaat menuliskan diagram terlebih dahulu:
diagram menjadi ukuran kebenaran program, bukan sekadar catatan.
\end{peringatanbox}
